\documentclass[10pt,a4paper]{article}

% ----------------- Paquetes -------------------%
\usepackage[T1]{fontenc}
\usepackage[utf8]{inputenc}
\usepackage[a4paper,margin=2.5cm]{geometry}
\usepackage{mathptmx}             
\usepackage{changepage} 
\usepackage{setspace}
\usepackage[spanish]{babel}
\usepackage[labelfont=bf,labelsep=space]{caption}
\usepackage{amssymb}
\usepackage{amsmath}
\usepackage{ebgaramond}
\usepackage{placeins}
\usepackage{etoolbox}
\usepackage{setspace}
\usepackage{parskip} 
\usepackage{tcolorbox}
\usepackage{needspace}
\usepackage{xparse}
\usepackage{ocgx2}
\usepackage{hyperref}
\usepackage{hypcap}
\usepackage{soul}\sethlcolor{yellow}% resaltado amarillo: \hl{texto} (Panel TCEE, ⌃H)


%%%%%%%%%%%%%%%%%%%%%%%%%%%%%%%%%%%%%%%%%%%%%%%%%%%%%%%%%%%%%%%%
% --- Fija primero tus valores globales "buenos" ---
\setstretch{1.2}                 % Interlineado global
\setlength{\parskip}{0.7\baselineskip}
\setlength{\parindent}{0pt}
\newlength{\GlobalParskip}
\setlength{\GlobalParskip}{\parskip}

\newcounter{eqnum}[subsection]
\renewcommand{\theeqnum}{\arabic{eqnum}}
\newcounter{alphaitem}
\newcounter{numitem}

%% ------------------- Recuadros----------------------%%
\newcommand{\puntosclave}[1]{%
  \begin{tcolorbox}[
    colframe=black,
    title={Puntos clave},
    before upper={\setlength{\parindent}{0.5cm}\setlength{\parskip}{0.5\baselineskip}}
  ]
  #1
  \end{tcolorbox}
}

\newcommand{\modificaciones}[1]{
  \begin{tcolorbox}[
    colback=white,
    colframe=red,
    title={Modificaciones a realizar},
    before upper={\setlength{\parindent}{0.5cm}\setlength{\parskip}{0.5\baselineskip}}
  ]
  #1
  \end{tcolorbox}
}

%% ------------------- Preguntas TEST----------------------%%
% Opciones: radio (single-choice)
\NewDocumentCommand{\OptRadio}{m m m}{%
  \noindent\CheckBox[radio,name=#1,value=#2,width=1.2ex,height=1.2ex]{}%
  \;\textbf{#2.}~#3\par
}

% Opciones: checkbox (multi-choice)
\NewDocumentCommand{\OptCheck}{m m}{%
  \noindent\CheckBox[name=#1,width=1.2ex,height=1.2ex,bordercolor={0 0 0}]{}%
  \;#2\par
}

% Pregunta single-choice (radio)
% Args: 1=id, 2=enunciado, 3=opciones, 4=solución+justificación, 5=imagen (o {}), 6=origen
\NewDocumentCommand{\PreguntaTestSingle}{m m m m m m}{%
  \par\medskip
  \noindent\textbf{#1.}~#2\par\smallskip

    % Imagen (si no es {})
    \IfBlankTF{#5}{}{%
      \begin{center}
        \includegraphics[width=0.75\linewidth]{#5}
      \end{center}
    }

    \begin{Form}
      #3
    \end{Form}

    \par\smallskip
    \switchocg{sol-#1}{\fbox{Mostrar / ocultar solución}}%
    \hfill{\footnotesize #6}\par\smallskip

    \begin{ocg}{Solución}{sol-#1}{0}
      \noindent #4
    \end{ocg}
  \par\medskip
}

% Pregunta multi-choice (checkboxes)
\NewDocumentCommand{\PreguntaTestMulti}{m m m m m m}{%
  \par\medskip
  \noindent\textbf{#1.}~#2\par\smallskip

    % Imagen (si no es {})
    \IfBlankTF{#5}{}{%
      \begin{center}
        \includegraphics[width=0.75\linewidth]{#5}
      \end{center}
    }
    \begin{Form}
      #3
    \end{Form}

    \par\smallskip
    \switchocg{sol-#1}{\fbox{Mostrar / ocultar solución}}%
    \hfill{\footnotesize #6}\par\smallskip

    \begin{ocg}{Solución}{sol-#1}{0}
      \noindent #4
    \end{ocg}
  \par\medskip
}

%% ------------------- CITA ----------------------%%
% \begin{cita}[Autor][Año][Obra] texto \end{cita}: cita sangrada y, debajo a la derecha, «AUTOR (Año), Obra». Los tres datos son opcionales.
% En el Panel TCEE: escribir \cita e Intro (el cursor va al texto; Tab pasa a Autor, Año y Obra).
\NewDocumentEnvironment{cita}{ O{} O{} O{} +b }{%
  \begin{adjustwidth}{2cm}{0pt}
  \raggedright
  #4\par
  \raggedleft
  \if\relax\detokenize{#1#2#3}\relax
    % sin firma
  \else
    #1%
    \if\relax\detokenize{#2}\relax\else\if\relax\detokenize{#1}\relax\else\space\fi(#2)\fi
    \if\relax\detokenize{#3}\relax\else\if\relax\detokenize{#1#2}\relax\else, \fi\textit{#3}\fi
  \fi
  \end{adjustwidth}
}{}



%% ------------------- Listas----------------------%%
    % --- Lista alfabética ---
    \newenvironment{la}
      {\begin{list}{\alph{alphaitem})}{%
          \usecounter{alphaitem}%
          \setlength{\leftmargin}{1cm}%
          \setlength{\parskip}{3pt}% 
          \setlength{\itemsep}{0pt}% ← espacio entre ítems
          \setlength{\parsep}{0pt}%  ← párrafos dentro de un ítem
      }}
      {\end{list}}
      
    % --- Lista numérica ---
    \newenvironment{lnum}
      {\begin{list}{\arabic{numitem}.}{%
          \usecounter{numitem}%
          \setlength{\leftmargin}{1cm}%
          \setlength{\parskip}{3pt}% 
          \setlength{\itemsep}{0pt}% ← espacio entre ítems
          \setlength{\parsep}{0pt}%  ← párrafos dentro de un ítem
      }}
      {\end{list}}
      
% ------------------- Imágenes----------------------%
    \graphicspath{{Img/}}% Rutas por defecto 
    % --- Imagen adaptativa ---
    % Registros de dimensión definidos una sola vez
    \newdimen\imgcap
    \newdimen\imgmin
    \newdimen\imgmax
    \newdimen\imgremain
    \newdimen\imgh
    \newdimen\imgneed
    
    \newcommand{\imagenfit}[3]{%
      \addvspace{10pt}% separación antes
      \par\begingroup
        % Parámetros de composición (asignaciones, no \newdimen)
        \imgcap=1\baselineskip            % alto estimado de título+fuente
        \imgmin=0.15\textheight
        \imgmax=0.15\textheight
        \imgremain=\pagegoal
        \ifdim\pagegoal=\maxdimen \imgremain=0pt\fi
        \advance\imgremain by -\pagetotal
        \advance\imgremain by -\imgcap
        % Decisión de altura destino
        \ifdim\imgremain<\imgmin
          \newpage
          \imgh=0.20\textheight
        \else
          \imgh=\imgremain
          \ifdim\imgh>\imgmax \imgh=\imgmax \fi
        \fi
        % Reservar espacio para evitar cortes del bloque completo
        \imgneed=\imgh \advance\imgneed by \imgcap
        \Needspace{\imgneed}
        % Cuerpo
        \noindent\begin{minipage}{\textwidth}
          \centering
          \captionof{figure}{\textemdash\ #2}\par\vspace{0.3em}% título arriba
          \includegraphics[width=\linewidth,height=\imgh,keepaspectratio]{\detokenize{#1}}\par
          \vspace{0.3em}\small\textit{Fuente: }#3% fuente abajo
        \end{minipage}%
      \endgroup
      \par\addvspace{10pt}%
    }

%------------ Bloque de RECUADRO ESPECIAL -------------%
    \newenvironment{specialblock}[1]{%
      \begin{quote} % crea sangría a izquierda y derecha
      \small % texto más pequeño
      \noindent\textbf{#1}\\[-0.3em] % título en negrita
      \rule{\linewidth}{0.4pt}\\[0.6em]}{
      \end{quote}}
      
%-------------------------- Portada -----------------------%
\newcommand{\oldtitle}[1]{\def\OldTitle{#1}}
\newcommand{\changerationale}[1]{\def\ChangeWhy{#1}}

\newcommand{\changebox}{%
\begin{tcolorbox}[title={Historial de cambios del tema}]
\textbf{Título anterior:} \OldTitle\par
\textbf{Motivo del cambio:} \ChangeWhy
\end{tcolorbox}
}

%-------------------------- Author Font -----------------------%
    \newcommand{\authorfont}[1]{{\fontfamily{EBGaramond-TLF}\selectfont #1}}

%-------------------------- Anexos -----------------------%
    \makeatletter
    \newcounter{anexo}
    \renewcommand{\theanexo}{\Roman{anexo}}
    \newcommand{\anexo}[1]{%
      \clearpage
      \phantomsection
      \refstepcounter{anexo}%
      \noindent\textbf{Anexo \theanexo.- }#1\par\medskip
      \addcontentsline{stc}{section}{Anexo \theanexo.- #1}%
    }
    \makeatother

    %-------------------------- Lista de figuras (personal) -----------------------%
\makeatletter
\newcommand{\MyListOfFigures}{%
  \clearpage
  \phantomsection
  {\centering\bfseries\Large Índice de figuras\par}%
  \medskip
  % Entrada en nuestro índice de contenido (stc)
  \addcontentsline{stc}{section}{Índice de figuras}%
  % Recuperación del listado (sin cabecera propia)
  \begingroup
    \parindent=0pt\parskip=2pt
    \@starttoc{lof}%
  \endgroup
}
\makeatother

%-------------------------- Bibliografía -----------------------%
\makeatletter
% Contador para las referencias
\newcounter{myref}
% Cabecera de la bibliografía, entra en el índice 'stc'
\newcommand{\MyBibliography}{%
  \clearpage
  \phantomsection
  {\centering\bfseries\Large Bibliografía y referencias\par}%
  \medskip
  \addcontentsline{stc}{section}{Bibliografía y referencias}%
  \setcounter{myref}{0}%
}
\newcommand{\MyBibItem}[4]{%
  \refstepcounter{myref}%
  \noindent[\themyref]\ #1.\ \emph{#2}. #3. #4\par
}
\makeatother

%--------- Establecer cuatro niveles de índice --------%
    \setcounter{tocdepth}{4}   
    \setcounter{secnumdepth}{4}% numera hasta \paragraph

%%%%%%%%%%%%%%%%%%%%%%%%%%%%%%%%

\makeatletter

    %--------- Interlineado Global --------%
    \edef\GlobalBaseLineStretch{\baselinestretch}
    
    %--------- Espaciado de seccinones --------%
    \renewcommand\section{\@startsection{section}
      {1}{\z@}
      {2.5ex \@plus 1ex \@minus .2ex} % espacio antes
      {1.5ex \@plus .2ex}             % espacio después
      {\normalfont\Large\bfseries}    % formato del título
    }
    \renewcommand\subsection{\@startsection{subsection}{2}{\z@}
      {3ex \@plus 0.8ex \@minus .2ex}
      {1.5ex \@plus .2ex}
      {\normalfont\large\bfseries}
    }
    \renewcommand\subsubsection{\@startsection{subsubsection}{3}{\z@}
      {3ex \@plus 0.5ex \@minus .2ex}
      {1ex \@plus .2ex}
      {\normalfont\normalsize\bfseries}
    }
    \renewcommand\paragraph{\@startsection{paragraph}{4}{\z@}%
    {-2ex \@plus -0.5ex \@minus -.2ex}% espacio antes
    {0.8ex \@plus .2ex}% espacio después
    {\normalfont\normalsize\itshape}}% ← cursiva en lugar de negrita

    %--------- Ecuaciones --------%   
    % Variables globales para almacenar títulos y notas
    \gdef\leq@current@section{}%
    \gdef\leq@current@subsection{}%
    \gdef\leq@last@printed@section{}%
    \gdef\leq@last@printed@subsection{}%
    
    % Comando para añadir nota (invisible en el cuerpo)
    \newcommand{\eqnote}[1]{%
      \addcontentsline{leq}{leqnote}{#1}%
    }
    
    % Comando para ecuaciones
    \newcommand{\eqblock}[2]{%
      % Verificar si necesitamos imprimir el título de sección
      \ifx\leq@current@section\leq@last@printed@section\else
        \addcontentsline{leq}{leqsection}{\leq@current@section}%
        \global\let\leq@last@printed@section\leq@current@section
        \gdef\leq@last@printed@subsection{}% Reset subsección al cambiar sección
      \fi
      % Verificar si necesitamos imprimir el título de subsección
      \ifx\leq@current@subsection\leq@last@printed@subsection\else
        \ifx\leq@current@subsection\@empty\else
          \addcontentsline{leq}{leqsubsection}{\leq@current@subsection}%
          \global\let\leq@last@printed@subsection\leq@current@subsection
        \fi
      \fi
      % Ecuación
      \refstepcounter{eqnum}%
      \begin{equation*}
        #1\tag{E.\theeqnum}
      \end{equation*}%
      \addcontentsline{leq}{leqt}{[E.\theeqnum] -- #2}%
      \addcontentsline{leq}{leqm}{#1}%
    }
    
    %%% ----- Índice de ecuaciones ----------%%%
    % Tamaño para []
    \newcommand{\leqIndexFont}{\normalsize}
    \newcommand{\leqTitleFont}{\tiny}
    \newcommand{\leqNoteFont}{\tiny}
    % Cabecera de sección 
    \newcommand*\l@leqsection[2]{%
      \addvspace{25pt}% Mayor separación antes de nueva sección
      \noindent\leqIndexFont
      \makebox[\linewidth]{\rule{.38\linewidth}{0.4pt}\ \ \textbf{  \MakeUppercase{#1}}\ \ \rule{.38\linewidth}{0.4pt}}%
      \par\vspace{-10pt}
    }
    % Cabecera de subsección 
    \newcommand*\l@leqsubsection[2]{%
      \par\addvspace{15pt}%
      \noindent\leqIndexFont #1
      \par\vspace{-7pt}
    }
    % Título de ecuación 
    \newcommand*\l@leqt[2]{%
    \par\addvspace{10pt}%
      \par\noindent\leqTitleFont #1\par
    }
    % Miniatura de ecuación
    \newcommand*\l@leqm[2]{%
      \vspace{0.3\baselineskip}%
      \par\scriptsize\noindent\hspace*{1.5em}\(#1\)\par%
    }
    % Nota de ecuación 
    \newcommand*\l@leqnote[2]{%
      \vspace{0.2\baselineskip}%
      \par\noindent\leqNoteFont\hspace*{1.5em}\textbf{Nota.--} #1\par\vspace{0.5\baselineskip}%
    }
    % Generación del índice
    \newcommand{\mylistofequations}{%
      \clearpage
      \phantomsection
      % Título visual de la lista (ajústalo a tu gusto):
      {\centering\bfseries\Large Índice de ecuaciones\par}
      % Entrada en nuestro índice 'stc':
      \addcontentsline{stc}{section}{Índice de ecuaciones}%
      % Llamada a tu lista real (usa el nombre exacto de tu macro):
      \@starttoc{leq}%
    }
    
    %--------- Captura de títulos de sección --------%
    \let\leq@original@section\section
    \let\leq@original@subsection\subsection
    % Flag para saber si estamos en una sección asterisco
    \newif\ifLeqStarSection
    \LeqStarSectionfalse
    
    % Redefinir \section CON CONTROL DE ASTERISCO
    \renewcommand{\section}{%
      \@ifstar{\leq@section@star}{\@dblarg\leq@section@nostar}%
    }
    \def\leq@section@star#1{%
      \global\LeqStarSectiontrue
      \gdef\leq@current@section{#1}%
      \gdef\leq@current@subsection{}%
      \leq@original@section*{#1}%
      \global\LeqStarSectionfalse
    }
    \def\leq@section@nostar[#1]#2{%
      \global\LeqStarSectionfalse
      \gdef\leq@current@section{#2}%
      \gdef\leq@current@subsection{}%
      \leq@original@section[#1]{#2}%
    }
    % Redefinir \subsection
    \renewcommand{\subsection}{%
      \@ifstar{\leq@subsection@star}{\@dblarg\leq@subsection@nostar}%
    }
    \def\leq@subsection@star#1{%
      \global\LeqStarSectiontrue
      \gdef\leq@current@subsection{#1}%
      \leq@original@subsection*{#1}%
      \global\LeqStarSectionfalse
    }
    \def\leq@subsection@nostar[#1]#2{%
      \global\LeqStarSectionfalse
      \gdef\leq@current@subsection{#2}%
      \leq@original@subsection[#1]{#2}%
    }

    % ----- Restauración de valores Globales de estilo ----------%
    % Flags
    \newif\ifInFootnote
    \InFootnotefalse
    \pretocmd{\@footnotetext}{\InFootnotetrue}{}{}
    \apptocmd{\@footnotetext}{\InFootnotefalse}{}{}
    % Profundidad de listas (soporta anidadas)
    \newcount\MyListDepth \MyListDepth=0
    \AtBeginEnvironment{la}{\advance\MyListDepth by 1}
    \AfterEndEnvironment{la}{\advance\MyListDepth by -1}
    \AtBeginEnvironment{lnum}{\advance\MyListDepth by 1}
    \AfterEndEnvironment{lnum}{\advance\MyListDepth by -1}
    \AtBeginEnvironment{itemize}{\advance\MyListDepth by 1}
    \AfterEndEnvironment{itemize}{\advance\MyListDepth by -1}
    \AtBeginEnvironment{enumerate}{\advance\MyListDepth by 1}
    \AfterEndEnvironment{enumerate}{\advance\MyListDepth by -1}
    % Restauración diferida: hazlo al INICIO del próximo párrafo real
    \newif\if@pendingRestore
    \@pendingRestorefalse
    \newcommand{\RequestRestore}{\global\@pendingRestoretrue}
    \everypar{%
      \if@pendingRestore
        \global\@pendingRestorefalse
        \ifInFootnote\else
          \ifnum\MyListDepth=0
            \setlength{\parskip}{\GlobalParskip}%
            \setstretch{\GlobalBaseLineStretch}%
          \fi
        \fi
      \fi
    }
    % Al final de bloques:
    \AfterEndEnvironment{equation}{\RequestRestore}
    \AfterEndEnvironment{equation*}{\RequestRestore}
    \AfterEndEnvironment{la}{\RequestRestore}
    \AfterEndEnvironment{lnum}{\RequestRestore}
    % Al final de macros:
    \apptocmd{\eqblock}{\RequestRestore}{}{}
    \apptocmd{\imagenfit}{\RequestRestore}{}{}
    
\makeatother

% ===================== ToC propio con 4 niveles (único bloque) =====================
\makeatletter

% --- Escritores: añadimos una línea al .stc en cada cabecera numerada ---
% Guardamos las versiones actuales para llamar al comportamiento normal
\let\MyTOC@orig@section\section
\let\MyTOC@orig@subsection\subsection
\let\MyTOC@orig@subsubsection\subsubsection
\let\MyTOC@orig@paragraph\paragraph

% SECTION
\renewcommand{\section}{%
  \@ifstar{\MyTOC@section@star}{\@dblarg\MyTOC@section@nostar}%
}
\def\MyTOC@section@star#1{%
  \MyTOC@orig@section*{#1}% sin entrada al .stc
}
\def\MyTOC@section@nostar[#1]#2{%
  \MyTOC@orig@section[{#1}]{#2}% comportamiento normal (escribe al .toc)
  % Escribimos también nuestra línea al .stc (texto con \numberline)
  \addcontentsline{stc}{section}{\protect\numberline{\thesection}#2}%
}

% SUBSECTION
\renewcommand{\subsection}{%
  \@ifstar{\MyTOC@subsection@star}{\@dblarg\MyTOC@subsection@nostar}%
}
\def\MyTOC@subsection@star#1{%
  \MyTOC@orig@subsection*{#1}%
}
\def\MyTOC@subsection@nostar[#1]#2{%
  \MyTOC@orig@subsection[{#1}]{#2}%
  \addcontentsline{stc}{subsection}{\protect\numberline{\thesubsection}#2}%
}

% SUBSUBSECTION
\renewcommand{\subsubsection}{%
  \@ifstar{\MyTOC@subsubsection@star}{\@dblarg\MyTOC@subsubsection@nostar}%
}
\def\MyTOC@subsubsection@star#1{%
  \MyTOC@orig@subsubsection*{#1}%
}
\def\MyTOC@subsubsection@nostar[#1]#2{%
  \MyTOC@orig@subsubsection[{#1}]{#2}%
  \addcontentsline{stc}{subsubsection}{\protect\numberline{\thesubsubsection}#2}%
}

% PARAGRAPH
\renewcommand{\paragraph}{%
  \@ifstar{\MyTOC@paragraph@star}{\@dblarg\MyTOC@paragraph@nostar}%
}
\def\MyTOC@paragraph@star#1{%
  \MyTOC@orig@paragraph*{#1}%
}
\def\MyTOC@paragraph@nostar[#1]#2{%
  \MyTOC@orig@paragraph[{#1}]{#2}%
  % Si no numeras los \paragraph, cambia \theparagraph por texto plano sin \numberline
  \addcontentsline{stc}{paragraph}{\protect\numberline{\theparagraph}#2}%
}

% --- Impresor del índice propio (lee el .stc) ---
\newcommand{\MyTOC}{%
  \par\bigskip
  {\centering\bfseries\Large Índice de contenido\par}%
  \medskip
  \begingroup
    \parindent=0pt\parskip=2pt
    \@starttoc{stc}%
  \endgroup
}

\makeatother
% ===================== fin ToC propio =====================


%%%%%%%%%%%%%%


% --- Datos del tema ---
\title{Análisis del crecimiento de la empresa\\
\large Métodos de valoración de empresas. Especial referencia a los procesos de fusión, adquisición y alianzas estratégicas}
\author{Marcelo Ortega}
\date{Marzo 2026}
\newcommand{\TemaCode}{3B4}

\oldtitle{
    B.4. Análisis del crecimiento de la empresa. Métodos de valoración de empresas. Especial referencia a los procesos de fusión, adquisición y alianzas estratégicas
}
\changerationale{
    Sin cambios.
}

\begin{document}


\maketitle
\changebox

\newpage

% --- Página de resumen y modificaciones ---
\puntosclave{ %% Escribir puntos clave Apuntes CECO
     La estructura del tema es clara, y consiste en dos bloques separados. En el primero, se ha optado por proveer al opositor de una visión conceptual y relativamente superficial, por lo que será recomendable profundizar (especialmente en lo referente a la operativización de la transacción): \textbf{no caer en una exposición que consista en una mera lista de definiciones} básicas: para garantizar el aprobado, será necesario un enfoque analítico.

     El segundo bloque trata de una tarea específica en toda operación de M\&A: determinar un precio para la adquisición. Es una cuestión técnica, y, como en el resto del tema, es relevante que la exposición alcance un correcto equilibrio entre los fundamentos teóricos puros y la práctica de la industria. Según el opositor incremente su conocimiento acerca de los contenidos de este título, seguramente incrementará también el peso que le dedicará a este bloque de la exposición: es de alguna manera el núcleo del tema, y donde se contienen los principales elementos diferenciales respecto a los 3B I a 3B3.

     Es recomendable que el opositor cuente con una actualización en cuanto a cifras y principales operaciones de fusiones y adquisiciones que se produzcan en el momento de estudiar el tema. Sin que resulte estrictamente necesario, permite dar un contexto ál contenido del tema.
    }
\vspace{1cm}

\modificaciones{
    Al introducir el primer apartado \textit{Análisis del crecimiento de la empresa}, es imprescindible añadir como motivaión: \textit{el culto a la reducción de costes}.

    \href{https://www.nber.org/system/files/chapters/c2059/c2059.pdf}{Muy interesante: Cifras desde los 60's (mercado US) + Anexo ICEX-CECO Cifras desde los 80's (concentración)}

    \href{https://postkeynesian.net/media/working-papers/PKWP2505.pdf}{Evidencia de la gran falacia. La creación de valor: Depende Inversión/Beneficios la composición del capital en las cuentas anuales} 
}


\newpage
\MyTOC
\newpage

\mylistofequations
\clearpage

\MyListOfFigures
\clearpage


\pagenumbering{arabic}
\setcounter{page}{1}


\section{Introducción}



\subsection{Contextualización}


\subsubsection{Relevancia}




\subsubsection{Problemática}



\subsection{Estructura de la exposición}






\clearpage
\section{Crecimiento inórganico: Modificaciones estructurales}
\puntosclave{
    El crecimiento inorgánico se puede definir en contraposición al orgánico: busca alcanzar el tamaño óptimo de la empresa a través de capacidades y recursos externos y mediante operaciones corporativas de fusión o adquisición.
    
    Las operaciones de M\&A se pueden clasificar como horizontales, verticales o conglomerados (según la relación entre los negocios) y como hostiles o amistosas (según la actitud de los equipos directivos). 
    
    Las fusiones y adquisiciones tienden a producirse en oleadas: influyen factores microeconómicos (shocks específicos) y macro-financieros (valoraciones del mercado bursátil). 
    
    Entre los principales determinantes microeconómicos encontramos la necesidad de crecer más rápido de lo que el crecimiento orgánico permite, las sinergias fruto de una mejora en la eficiencia corporativa y los motivos de agencia (incentivos personales de los directivos) 
    
    Las operaciones de fusión y adquisición entre compañías se pueden estructurar a través de un canje de acciones o mediante un pago en efectivo (con la consiguiente necesidad de financiación). En la elección entre ambas alternativas influyen elementos corno las valoraciones bursátiles o el reparto de riesgo y rentabilidad entra los accionistas de la empresa adquirida y la adquirente 
    
    La cooperación y las alianzas estratégicas se constituyen como una alternativa al crecimiento externo a través del M\&A. Pemliten una mayor flexibilidad para la cooperación empresarial pero a su vez pueden enfrentar problemas de gobernanza
}

Para acompañar adecuadamente este tema, lo mejor es imaginarse que somos el gerente de una empresa. Como gerentes tendremos diversos objetivos, pero, indudablemente, uno de ellos será o estará de forma indirecta relacionado con aumentar el valor de la empresa:

\eqblock{
\begin{aligned}
    &\text{OBJETIVO}: \text{Crecimiento}\equiv g\quad \Longrightarrow\quad \underset{\text{\scriptsize Las empresas se financiancian para \textbf{crecer}}}{\boxed{g=\underbrace{\left[\text{ROA}+(\text{ROA}-i)\frac{D}{E}\right]}_{= \text{ROE}}(1-d)}}
\end{aligned}
}{Objetivo: Crecimiento de la empresa y descomposición. Fuentes de financiación}

Una de las formas que tenemos para aumentar el valor de nuestra empresa es a través de la inversión en nuestro propio ecosistema. Es decir, llevar a cabo operaciones que aumenten el tamaño de nuestra empresa a través de la expansión de sus operaciones existentes. Es lo que se conoce como crecimiento orgánico.

En contraposición, el crecimiento inorgánico busca alcanzar los mismos objetivos pero a través de la adquisición completa de otras empresas, la adquisición de divisiones de negocio y/o de activos específicos. Los motivos que llevan a optar por el crecimiento inorgánico (externo) son variados como veremos después, pero la diferencia clave con el crecimiento interno es la \textbf{velocidad en su ejecución}.

En este tipo de operaciones, por tanto, no nos interesará tanto el rendimiento esperado de la inversión (puesto que responde al plan estratégico individualizado de cada empresa) sino la valoración de la empresa objetivo. Pues será determinante para asegurar el rendimiento de la operación y, en todo caso, escoger entre alternativas de financiación.\footnote{%
    En el lenguaje de las normas internacionales de contabilidad (NIIF). se denominan combinaciones de negocios y de manera genérica en el ámbito del análisis financiero, se denomina actividad de M\&A (Mergers \& Acquisitions).
}

\textbf{¿Qué queremos decir?} -- Por tanto, sabiendo cuál es nuestro ojetivo como gerentes, y el objeto del análisis que emprendemos (el crecimiento inorgánico), lo primero que debemos saber es las opciones de las que disponemos, tanto desde un punto de vista jurídico como económico. Esto es esencial puesto que conocer el tipo de operaciones y sus particularidades permite tomar mejores decisiones atendiendo a las circunstancias que nos ocupan.

\subsection{Operaciones con alteración Patrimonial: Fusiones, Adquisiciones y Reestructuraciones}
El derecho mercantil y fiscal en las distintas jurisdicciones tiene una influencia clave en los negocios jurídicos que se utilizan para estructurar estas operaciones. Sin embargo, optamos por un enfoque estríctamente económico-financiero, más adecuado en nuestro caso. Estríctamente hablando, crecer externamente puede realizarse por dos vías: a través de la fusión con otras empresas y a través de la adquisición de estas.

\subsubsection{Fusiones: Consolidación de Patrimonios}
Una fusión es una operacion de concentración por la cual dos sociedades independientes consolidan sus patrimonios y actividades bajo una sola entidad. Se trata de operaciones en las que habitualmente el tamaño de ambas compañías es similar o al menos muy significativo en términos relativos de una frente a la otra. 

En casi todas las jurisdicciones hay dos clases posibles de fusiones (incluido en el caso de España):
\begin{la}
    \item \textbf{Fusión por creación}. En la que nace una nueva empresa con personalidad jurídica propia extinguiéndose las de las anteriores; y,
    \item \textbf{Fusión por absorción}. En la que se mantiene la personalidad jurídica de la absorbente y se extingue la de la absorbida. 
\end{la}

\subsubsection{Adquisiciones: Compra Patrimonial}  
Por otro ldo, las adquisiciones concisten en la compra de todos o parte de los activos de otra compañía. Por tanto, se suele hablar de adquisición cuando el tamaño es clarámente desigual: una compañía grande adquiere a una más pequeña.\footnote{%
    En ocasiones, la diferencia con una fusión por absorción puede no estar clara, pero se trata de diferencias jurídicas en la forma de la transacción. En el caso español por ejemplo, mientras que las fusiones por absorción se denominan con ese nombre, la adquisición, entendida como la compra de todo el patrimonio de una compaía por otra, se regula en la figura de la cesión global de activo y pasivo.
} De acuerdo con la actitud de los equipos directivos las operaciones pueden categorizadas como:
\begin{la}
    \item \textbf{Carácter hostil}. Intentos de adquirir una empresa en contra de la opinión de sus directivos. Esta reacción puede deberse al precio ofertado o a razones de otra índole (en ocasiones vinculadas al futuro personal de los. propios empleados);
    \item \textbf{Carácter amistoso}. Cuando cuenta con el respaldo de los directivos de ambas empresas. No obstante, esto no es una garantía de que la operación vaya a llevarse a cabo, dado que los accionistas pueden no estar de acuerdo con las condiciones. ofertadas o con la operación en términos generales.
\end{la}

\paragraph{Empresas cotizadas: Oferta Pública de Adquisición}
Cuando la compañía adquirida es una empresa cotizada, el adquirente tiene que llevar a cabo una oferta pública de adquisición (OPA), cumpliendo con los requisitos que establecen sobre estas ofertas las distintas legislaciones mercantiles, de competencia y supervisión del mercado. 

En las cotizadas, las principales diferencias entre la fusión y la adquisición radica en que en una adquisición los accionistas pueden ser apelados directamente por la empresa adquirente, sin necesidad del consentimiento de los directivos de la empresa objetivo y siempre de acuerdo con la legislación correspondiente. En las fusiones sin embargo, primero hay una negociación entre los directivos y una vez alcanzado el acuerdo este debe ser confirmado por los accionistas.

\paragraph{Conglomerados empresariales: \textit{Holdings}}
Un tipo particular de operación de adquisición que puede resultar en patrimonios diferenciados son los conglomerados empresariales o holdings. Es cada vez más recuente encontrarnos con empresas que pertenecen a una empresa matriz cuyo objeto social está tremendamente diseminado: una misma empresa es propietaria de una serie de empresas con procesos productivos distintos y no (necesariamente) relacionados ni con el negocio principal de la compañía ni con el negocio de las subempresas que la conforman. 

Las operaciones financieras de un holding son dificilmente encuadrables dentro del marco tradicional de la Teoría de la Organización Indutrial o de la empresa ya que este holding, que puede llevar a cabo las operaciones con o sin ánimo de control, persigue más bien el objetivo de tenencia de una cartera de empresas, no de una actividad productiva concreta.\footnote{%
    Las operaciones en las que entra un conglomerado o holding suelen ser más baratas que las fusiones o adquisiciones porque no suele existir una prima de control y no es necesario hacerse con la mayoría o totalidad del capital. El sentido económico de estos conglomerados en ocasiones es puesto en duda, al tratarse de instrumentos de concentración de poder político-económico que no necesariamente responden a razones de eficiencia financiera y operativa. Por lo general, cuando una empresa con un "core business" bien definido, comienza a adquirir paquetes de acciones y en el mercado se percibe que se está convirtiendo en un holding, pasa a cotizar con un descuento por la percepción de que el inversor es más eficiente diversificando.
} 


\subsubsection{Crecimiento negativo: Desinversión}
Los procesos de crecimiento empresarial pueden tener ambas direcciones y en ocasiones las empresas buscan reducir su tamaño, dando lugar a procesos de desinversión. 

Estos procesos pueden ser necesarios por diversas razones, como cambios en la estrategia corporativa, necesidades de liquidez, reestructuración o eliminación de divisiones no rentables, entre otros. En estas circunstancias, las empresas pueden decidir eliminar operaciones completas o cerrar divisiones específicas que no son viables, lo que podría implicar la liquidación de activos y costosos procesos de destrucción de empleo (liquidación o restructuración). Por ello, en estos casos es mucho más habitual recurrir al (de)crecimiento externo, por la facilidad para deshacerse de divisiones de negocio completas mediante \textbf{operaciones corporativas}. 

Además de procesos de venta, también se recurre a \textit{spin-offs} (escisiones de activo y pasivo). Además de ser un medio para alcanzar el tamaño óptimo empresarial y reasignar capital, estas operaciones de desinversión también pueden ser una fuente de información para el conjunto del mercado. Por un lado, permiten poner en valor activos y divisiones de negocio rentables, y por otro lado, pueden ayudar a reducir los problemas de información asimétrica que generan

\begin{specialblock}{Escisiones y violación de la Hipótesis de los Mercados Eficientes: ¿la ley del precio único existe?}
    Un interesante suceso ocurrió precisamente a raíz de una operación de desinversión. 

    [Extraído de THALER (2001)] Una empresa tradicional de fabricación de suministros de red se fusionó con una empresa de fabricación de PDA's, ayudantes digitales (en concreto 3Com y Palm). Durante el verano de 1999, cuando el valor de las acciones de cualquier empresa tecnológica estaban por las nubes, las acciones de la empresa matriz estaban estancadas. Para solucionar este hecho, los directores de la empresa decidieron segregar la empresa con un plan muy concreto: se sacaría al mercado únicamente el 5\% de la compañía y, pasados unos meses, cada accionista de la empresa matriz recibiría 1,5 acciones de la empresa escindida por cada acción de la empresa matriz: ¿qué diferencia supondría esto? Una diferencia colosal.
    
    En un primer momento, al anunciar el plan de escisión, las acciones de la empresa matriz subieron un 150\%. Al sacar al mercado (el 5\%) la empresa escindida por 38\$, los precios de las acciones de la escisión subieron rápidamente hasta equipararse con la empresa matriz (en 100\$), que a su vez bajó ese mismo hasta situarse en 83\$: ¿qué esta pasando? ¿qué coherencia tiene esto? El mercado estaba valorando la empresa matriz con unas pérdidas de 23.000M\$, un negocio perfectamente rentable y mucho más maduro que el de la empresa matriz.

    \imagenfit{PALM_3COM.jpg}{La peculiar aritmética de PALM y 3COM}{\textit{Todo lo que he aprendido con la psicología económica}. THALER (2001)}
    
\end{specialblock}

\subsection{Operaciones sin alteración Patrimonial: Alianzas y cooperaciones}
También podemos observar operaciones sin alteración patrimonial que se llevan a cabo con el objetivo de aumentar el crecimiento de la empresa: las alianzas estratégicas o las cooperaciones. Al no haber alteraciones en el Patrimonio este se mantiene claramente diferenciado. 

Se trata de una parada intermedia entre crecimiento orgánico e inorgánico, en la que dos o más empresas deciden actuar conjuntamente para lograr determinados objetivos, compartiendo recursos, capacidades o actividades. Sin ánimo de ser exhaustivos, las razones que conducen a estas situaciones suelen tener que ver con: (i) cambios tecnológicos; (ii) acceso (recíproco, por ejemplo) a mercados tanto nacionales como internacionales; y, (iii) financiación compartida de grandes proyectos de inversión (UTE's, en España). 


Los acuerdos de cooperación pueden adoptar una personalidad jurídica propia que formaliza el compromiso entre los cooperantes: 
• Consorcio: Alianza entre varias empresas que mediante un contrato formaliza una relación a largo plazo entre cada una de las empresas y una organización conjuntaintegrada por ellas mismas. El objetivo es llevar a cabo una gran inversión que no podría efectuarse por las dos partes por separado.
• Joint-Venture: Dos o más empresas crean una tercera empresa independiente, concapital conjunto. Otra forma muy habitual de sellar estas alianzas estratégicas es mediante la adquisición cruzada de participaciones accionariales. La forma de estructurar estos acuerdos es variada pero suele implicar algún tipo de pacto en el que se refleje la gobernanza conjunta y los acuerdos de colaboración a nivel económico, técnico y financiero.

\begin{specialblock}{}
    Son muchos los ejemplo que podríamos poner en este campo, pero uno particularmente importante (e ilustrativo) se encuentra en el mundo de la aviación. 

    [Extraído de GAY DE LIEBANA, 2007] Son conocidos los corsés jurídicos sobre posibles fusiones entre compañías aéreas, donde de un lado existen normas restrictivas del transporte aéreo internacional que ponen barreras para que se lleven a cabo fusiones, integraciones y adquisiciones transfronterizas, a lo que se tiene que añadir el hecho de que sean bastantes los países que restringen la propiedad extranjera de las líneas aéreas. A pesar de ello el secreto para aumentar los ingresos y reducir los costes es muy sencillo: alianzas. Así, hoy, Oneworld donde participa la compañía de bandera española Iberia, entre otras, compite con SkyTeam, con Air France y Alitalia entre otras compañías, y Star Alliance, que cuenta, también entre otras, con la alemana Lufthansa, SAS y la española Spanair. ¿Las claves de esas tres grandes alianzas del tráfico áereo? El número de empleados -260.500, 186.000 y 284.475, respectivamente -, los destinos en total - 570, 512 y 700, respectivamente - el control de aeropuertos que constituyen referentes en el mundo de la aviación - Heathrow, Zurich, Madrid - y son hub, descollantes centros de distribución, con lo cual se planta cara al auge de las aerolíneas de vuelos baratos, a la vez que estas integraciones no causan la pérdida de "slots", es decir, de derechos de vuelo que suponen uno de los intangibles más apetecibles en el sector de la aviación comercial. El hecho de poder compartir estructuras, infraestructuras, servicios de handling, catering..., entre las compañías integradas en una misma alianza implica una mejora económica sustancial, aprovechando vuelos y recorridos así como las centrales de reservas.

    
\end{specialblock}



\paragraph{Ventajas: Riesgo y autonomía}
Por lo que respecta a las ventajas que estas operaciones suponen frente a las convencionales adquisiciones y/o fusiones son:
\begin{lnum}
    \item \textbf{Compartir riesgos}. Al cooperar o formar alianzas estratégicas, las empresas pueden compartir los riesgos y costos asociados con nuevas iniciativas o proyectos conjuntos;
    \item \textbf{Recursos}. Mediante estas operaciones, las empresas pueden acceder a los recursos, capacidades y conocimientos técnicos de sus socios, lo que les permite ampliar su alcance y competir de maneramás efectiva en el mercado;
    \item \textbf{Sinergias}. Estas operaciones pueden conducir a la creación de nuevos productos o servicios innovadores y a la expansión de su base de clientes; y,
    \item \textbf{Flexibilidad y autonomía}. Sobretodo, a diferencia de las fusiones y adquisiciones, permiten a las empresas mantener su independencia y autonomía jurídico-financiera, al conservar su identidad y estructura organizativa, mientras colaboran en áreas específicas de interés común.
\end{lnum}


\paragraph{Inconvenientes: }
Por lo que respecta a los inconvenientes, estas operaciones suelen tener un impacto desigual para los socios. Además, como cualquier otra interacción social, estas sujetas a riesgos inherentes a la conducta, como por ejemplo:
\begin{lnum}
    \item \textbf{Complejidad de la colaboración}. La colaboración entre empresas implica la gestión de relaciones complejas. Las diferencias en las culturas organizativas, los objetivos estratégicos y las prácticas comerciales pueden dificultar la cooperación efectiva y requerir una atención constante que puede drenar demasiados recursos de ambas organizaciones;
    \item \textbf{Rendimiento y expectativas}. En muchas ocasiones el éxito depende del desempeño y compromiso de los socios involucrados. Si uno de los socios no cumple con sus responsabilidades o no aporta valor esperado, puede afectar negativamente la alianza en su conjunto; y,
    \item \textbf{Compromisos}. Llegar a un acuerdo mutuamente beneficioso puede ser un desafio en las alianzas. Las empresas deben negociar y acordar términos y condiciones que satisfagan a ambas partes, lo que puede llevar tiempo y esfuerzo. Además, los intereses pueden cambiar a lo largo del tiempo, generando una tensión para diverger o diferenciarse de nuevo.
\end{lnum}

\subsection{Clasificaciones accesorias}
Además de las dos grandes categorías presentadas, existen otras formas de clasificar las operaciones de crecimiento externo. Sobre este aspecto, presentaremos principalmente dos clasificaciones: por las formas de adquisición y por el sector operativo de la empresa adquirida.

\subsubsection{Cadena de valor: Relación de negocios}
Desde una óptica de Política Económica y la valoración de las consecuencias de una fusión es importante tener en cuenta la posición relativa en la cadena de valor.

\paragraph{Intrasectoriales: Horizontal}
De acuerdo con esta relación, una operación se considera intrasectorial, u de integración horizontal, si se realiza entre empresas que pertenecen al mismo proceso productivo, habitualmente competidores. Las razones que llevan a este tipo de operaciones se analizan posteriormente en los determinantes del crecimiento inorgánico. 

\paragraph{Intersectoriales: Vertical}
Si la operación se lleva a cabo entre empresas con un proceso productivo diferente, u objeto social diferente, hablamos de integración vertical u intersectorial (a lo largo de la cadena de valor). Por ejemplo, un proveedor o un distribuidor. 

De esta forma, si el adquirente compra un objetivo que está por delante de él en la cadena de valor (un proveedor), se denomina integración hacia atrás (backward) y cuando un adquirente compra una empresa que está más abajo en la cadena de valor (un distribuidor), se denomina integración hacia delante (forward).

\subsubsection{Métodos de pago y financiación: Operativización}
Además, las operaciones pueden financiarse con efectivo, valores o una combinación de ambos en lo que se denomina una oferta mixta.

\paragraph{Oferta de efectivo}
En una oferta en efectivo, el efectivo puede proceder de los activos existentes de la empresa adquirente o de una emisión de deuda. De esta forma, el principal método para realizar una operación en efectivo es mediante préstamos puente, considerando como estos:\footnote{%
    Se trata de un préstamo temporal cuyo objeto es cubrir la necesidad inminente de financiación, pero sin una vocación de pemianencia, siendo necesario la búsqueda de una financiación a largo plazo definitiva. Permite conseguir los fondos con agilidad pero requiere una refinanciación.
}
\begin{lnum}
    \item \textbf{Ampliación de capital}. La empresa adquirente realiza una emisión de nuevas acciones. Es la opción con más riesgo para los bancos que participan en el préstamo puente y es habitual un sondeo previo del interés por participar en dicha ampliación;
    \item \textbf{Financiación a largo plazo}. En el menor plazo posible se estructura una financiación normalmente sindicada, que permita amortizar el préstamo puente. Los bancos que participan en la financiación sindicada a largo no tienen porque coincidir con los del préstamo puente. Como alternativa a la financiación bancaria la empresa puede emitir bonos, pero esto suele ser menos frecuente en operaciones de gran volumen; o,
    \item \textbf{Venta de activos}. Es un mecanismo cada vez más habitual por el que parte del préstamo puente se repaga con la venta de activos o filiales no consideradas estratégicas por la nueva dirección. También se puede estructurar mediante spin-offs que sacan a bolsa o venden parcialmente a terceros una participación en negocios estratégicos. Ha sido un mecanismo recurrente en el sector concesionario y constructor español con sus filiales de energías renovables.
\end{lnum}

\paragraph{Leverage Buy Out}
Además, es necesario destacar el papel cada vez más importante de las operaciones por compra apalancada. 

Muchas operaciones corporativas están protagonizadas por inversores financieros y fondos de capital riesgo, que aportan capital e impulsan cambios en la estrategia y dirección de las compañías. El mecanismo habitual al que recurren este tipo de agentes para financiar las fusiones y adquisiciones es el Leveraged Buyout (LBO) o compra apalancada. El objetivo de esta estructura es que la adquisición se financie endeudando (apalancando) a la propia compañía comprada, de tal modo que sean los propios recursos generados por la compañía comprada los que permitan repagar la financiación. Es frecuente que, una vez repagada la mayoría de la deuda, el inversor busque salir de la inversión (vendiendo la compañía o sacándola la bolsa, entre otras estrategias de salida). Al margen de la estrategia corporativa del adquirente, en última instancia se busca maximizar la rentabilidad del capital del inversor financiero.



\clearpage
\section{Valoración de empresas: Métodos}
\puntosclave{
    La valoración de las empresas es un elemento clave en los procesos de negociación en las fusiones y adquisiciones. Se deben tener en cuenta las sinergias esperadas en esta valoración. 
    
    El Descuento de Flujos de Caja es el método más utilizado, dado que permite obtener el valor intrínseco de una compañía de manera directa. Requiere la elaboración de un modelo financiero para la proyección de los flujos de caja futuro, que deberán ser descontados a la tasa de descuento correspondiente (el WACC). 
    
    A través del descuento de los flujos de caja libre (caja disponible tras hacer frente a las necesidades de inversión) se obtiene el Enterprise Value (el valor conjunto de la deuda y patrimonio neto de la compañía). Este es el valor de referencia en las operaciones corporativas.
    
    También existen métodos de valoración relativo, como la valoración por múltiplos. A partir de un universo de empresas o de transacciones comparables, se pueden obtener los múltiplos cotizados o pagados frente a alguna variable como el EBITDA, la cifra de ventas o el beneficio neto. Aplicar estos múltiplos a las compañías analizadas permite obtener una valoración indirecta.
    
    Las opciones reales complementan los métodos de valoración al integrar el valor de la flexibilidad en las decisiones empresariales.
}

Ahora que conocemos la arena conceptual en la que nos movemos y los tipos de operaciones que podemos llevar a cabo, tanto desde una perspectiva jurídico formal como desde una perspectiva económico-financiera, la duda inmediata es: ¿cómo realizarlas?

Volvamos al ejemplo inicial. Supongamos que somos el gerente de una empresa de almacenaje. En particular, nuestro objeto social es la creación de soluciones de almacenaje clásico y tenemos una buena cuota de mercado, sin restricciones significativas de financiación. Le hemos echado el ojo a una pequeña (pero pontente) empresa de logística que está trabajando en robots industriales. Como gerentes, sabemos que adquirir esta tecnología puede añadir mucho valor a nuestra empresa, sobretodo si, continuando el desarrollo, expandimos nuestras lineas de negocio ofreciendo automatismo en la gestión de los almacenes que instalamos. Pero, ¿cómo podemos valorar ese \textbf{beneficio portencial}? Es decir, ¿cómo podemos valorar la operación de forma que sea rentable tanto para nosotros como para la empresa?\footnote{%
    ¿Por qué tiene que ser mutuamente beneficiosa? Esta es la principal diferencia entre la valoración de una compañía como activo de inversión y su valoración a efectos de una operación corporativa.

La valoración de la empresa objetivo es el punto de partida de la negociación, de vital importancia para los analistas y accionistas de ambas partes, ya que determina el valor incial de la oferta. En este valor se tienen en cuenta aspectos como la distribución de beneficios en una fusión, basándose en las sinergias esperadas, o la \textbf{prima de control} a pagar por encima del valor intrínseco de la empresa objetivo. Por tanto, esta prima de adquisición debe reflejar las sinergias, reducciones de coste esperadas y la recompensa a los accionistas actuales por ceder el control. En esta segunda fase entran en juego múltiples aspectos, más allá de la pura valoración, dentro del proceso de negociación entre ambas partes para que la oferta sea equitativa y adecuada.
}

\textbf{¿Qué queremos decir?} A continuación presentaremos diferentes métodos que podemos utilizar para valorar la operación. No obstante, cabe advertir que se presentarán de forma genérica y, en la gran mayoría de casos, no será una cuestión de sustituir parámetros por valores sino que requerirán un estudio previo y ajustado al sector concreto y operación que enfrentemos. Además, en todo momento debemos tener presente que un proyecto de inversión no es equiparable a una operación de este tipo puesto que en ella rigen los \textbf{principios de la negociación}, de forma que un elemento crucial a considerar será la \textbf{prima de control}.

En términos jurídicos, las operaciones y su ejecución son diferentes si se trata de una empresa cotizada o no. De igual forma, los métodos de valoración serán diferentes. 

\subsection{Empresas no cotizadas}
Empecemos primero por los métodos de valoración característicos de las empresas no cotizadas. El punto de partida conceptual es muy simple: queremos valorar la capacidad del negocio de generar caja en condiciones normales de operación. Esto exige que nos distanciemos de cómo está la sociedad en su balance y optemos por valoraciones económicas que reflejen mejor su situación. 

Lo más habitual en este sentio es buscar el valor de la empresa, descontando su estructura de financiación y las desviaciones puntuales de circulante (fondo de maniobra), llegando a obtener el \textbf{enterprise value}.\footnote{%
    En el ámbito de las fusiones y adquisiciones el concepto sobre el que pivotan todos los métodos de valoración de la empresa a adquirir es el enterprise value. Es el valor que tiene la empresa para todos sus proveedores de capital: acreedores financieros y accionistas. Otras definiciones como el valor contable o el de liquidación son siempre accesorias y tienen un carácter secundario frente a este. Los acreedores financieros proveen fondos ajenos con un coste explícito y determinado contractualmente, habitualmente en la forma de deuda financiera. Por su parte, los accionistas suministran los fondos propios (capital social y beneficios no distribuidos) y su rentabilidad procede de los flujos residuales que no corresponden a los acreedores financieros.
} En resumidas cuentas:

\eqblock{
\begin{aligned}
    \boxed{\underset{\text{\scriptsize valor para los accionistas}}{\text{Equity Value}}=\underset{\text{\scriptsize valor para todos los acreedores financieros}}{\text{Enterprise Value}}-\text{NetDebt (Ajustes)}}
\end{aligned}
}{Valoración de la operación: Empresas no cotizadas}

Ese valor representa cuánto vale la actividad ordinaria si la aislamos de la deuda, la caja y de si en ese momento hay más o menos clientes por cobrar o proveedores por pagar de lo habitual. Pero, ¿por qué? En esencia, el valor de una empresa corresponde al valor que tiene la empresa para todos sus proveedores de capital: acreedores financieros y accionistas. Por eso, determinar este valor es esencial para poder ralizar posteriormente los ajustes necesarios que acaben proporcionando el valor de la empresa para nuestro vendedor objetivo: los \textbf{accionistas}.
Así que surge una pregunta esencial: ¿cómo podemos calcular el \textbf{enterprise value}? Existen diferentes métodos de valoración así que, sin ánimo de ser axhaustivos, presentaremos los más habituales.

\subsubsection{Descuento de Flujos de Caja}
El primero y más habitual de todos es la valoración por Descuento de Flujos de Caja (DFC). En esencia, se trata de obtener el enterprise value como un el \textbf{valor actual neto de los flujos monetarios futuros}: traer a valor presente flujos futuros. De esta forma, el enterprise value será el valor estimado de la corriente de flujos de ingresos netos futuros, que se atribuiría al conjunto de los acredores financieros. Medimos estos flujos a través de los denominados flujos de caja libre (o free cash flow , FCF en adelante).\footnote{%
    El flujo de caja libre (FCF) es la variable pertinente en este contexto porque representa el flujo de caja real que estaría a disposición de los inversores de la empresa después de realizar todas las inversiones necesarias tanto para acelerar el crecimiento de la empresa como para mantener la empresa operando a lo largo del tiempo con su dimensión actual. Estrictamente hablaríamos del Free Cash Flow to the Firm, el flujo de caja libre que permite analizar el valor de la compañía con independencia de su estructura de capital.

Además, en la terminología habitual de M\&A se conoce esta intertemporalidad como growth capex y maintenance capex (capital expenditure).
} 

Hay dos grandes etapas en un proceso de valoración: (i) analizar el modelo y plan de negocio; y, (ii) construir el modelo financiero de valoración. 

\paragraph{Primera etapa: Análisis del plan de negocio}
En la primera etapa se analizan los factores clave en la evolución del negocio para poder construir el conjunto de suposiciones e hipótesis que permitirán proyectar los ingresos, gastos, inversiones y estructura financiera. 

Hay dos elementos clave en esta etapa:
\begin{lnum}
    \item \textbf{Información de la empresa objetivo}. Es un trabajo que muy frecuentemente se contrata a consultores y especialistas externos para poder llevar a cabo las denominadas \textit{Due Diligence}. Estas evaluaciones abarcan numerosos aspectos desde lo comercial a lo fiscal, laboral, contable o medioambiental. El objetivo es garantizar la calidad de la información y reducir las asimetrías informacionales ante posibles contingencias o pasivos ocultos; y,
    \item \textbf{Análisis de negocio y sectorial}. Es aquí donde más importancia tiene la experiencia, intuición y capacidad del analista para dotar de credibilidad y calidad al análisis. Los aspectos técnicos y regulatorios juegan un papel muy importante en el análisis de las barreras de entrada y las ventajas competitivas.
\end{lnum}

\paragraph{Segunda etapa: Construcción del modelo financiero}
La segunda etapa es la elaboración del modelo financiero, que a su vez suele estructurarse en dos niveles: un primer nivel en el que se proyectan los flujos hasta un determinado horizonte temporal (normalmente, cinco años); y un segundo nivel en el que se estima el valor terminal o residual de los flujos de caja posteriores al primer nivel. Estas estimaciones de flujos de caja libres se descuentan a su valor actual y la suma de las dos partes es el valor estimado de la empresa.\footnote{%
    El arte del análisis financiero implica la capacidad de utilizar las herramientas adecuadas y ejercer el buen juicio para producir las mejores estimaciones posibles para cada partida de los estados financieros. En el proceso, los analistas hacen ajustes a sus proyecciones previas basándose en las sinergias propuestas y en los planes anunciados para la empresa fusionada. Por ejemplo, la duplicación de recursos podría dar lugar a la venta de una de las divisiones de la empresa objetivo. O los costes de explotación podrían ajustarse a la baja en previsión de economías de escala. Estos aj ustes son más fáciles de estimar en las fusiones amistosas, en las que el analista tiene acceso a datos financieros detallados sobre la empresa objetivo, que en las fusiones hostiles.
}\textsuperscript{,}\footnote{%
    En el análisis de empresas son particularmente importantes las previsiones de los niveles de precios (inflación, y tipos de interés futuros). Los ingresos se proyectan desglosando por tipo de producto, canal de venta o zona de geográfica y de la misma manera se realiza con los costes variables. Los gastos generales se suelen estimar por separada. Las necesidades de inversión son de dos tipos: (i) inversión en inmovilizado, tanto de expansión como de mantenimiento (denominado CAPEX); e, (ii) inversión en circulante. El capex de mantenimiento está vinculado a la reposición del inmovilizado que se va amortizando. Por su parte la inversión en circulante se estima a partir de la proyección del balance, tomando las diferencias entre pagos a proveedores, inventarios, cuentas a cobrar (es decir las partidas del activo y pasivo corriente). Con estas proyecciones de resultados y de balance se corrigen las partidas para llegar al flujo de caja libre.
}

\eqblock{
\begin{aligned}
    &\text{BN}_0\to\text{BN}_T\\[0.5em]
    &\begin{aligned}
        &(+)&&i_0,\ldots&&\to i_T
    \end{aligned}\\[0.5em]
    &\underset{\substack{\text{\scriptsize resultado operativo}\\\text{\scriptsize no apalacandado}}}{\text{NOPAT}_0}\to\text{NOPAT}_T\\
    &\over\begin{aligned}
        \\
        &(+)&&\text{Amortizaciones}_0,\ldots&&\to \text{Amortizaciones}_T\\[0.5em]
        &(-)&&\text{Fondo Maniobra}_0,\ldots&&\to \text{Fondo Maniobra}_T\\[0.5em]
        &(-)&&\text{CAPEX}_0,\ldots&&\to \text{CAPEX}_T
    \end{aligned}\\[0.5em]
    &\over\begin{aligned}
        \\
        &\text{Flujos de Caja Libre (FCL)}_0\to\text{Flujos de Caja Libre (FCL)}_T\\
    \end{aligned}\\[0.5em]
\end{aligned}
}{Proyección: Flujos de Caja Libres (CFL). Valoración de empresas: No cotizadas}

El método es sencillo, una vez obtenemos las proyecciones del paso previo, construimos la simulación financiera de la empresa. En esta simulación deberemos tener en cuenta: (i) los resultados de explotación proyectados; (ii) los gastos financieros (para no duplicar el escudo fiscal, que se considera en el coste del capital); (iii) las amortizaciones; (iv) las variaciones ajustadas del fondo de maniobra (working capital, inversión en circulante); (v) el CAPEX de expansión y mantenimiento y el cambio en el fondo de maniobra ; y, (vi) obtenedremos así el Flujo de Caja Libre (FCL).\footnote{%
    Este flujo de caja libre que no tiene en cuenta el beneficio fiscal de la deuda puede también denominarse Free Cash Flow to the Firm o Unlevered Free Cash Flow Es decir el flujo de caJa dcsapalancado.
}

Para descontar los Flujos de Caja deberemos recurrir a la estimación del coste del capital promedio (WACC). Para ello, existen diversas opciones, pero, sin ánimo de entrar a valorarlas consideraremos que es el tipo de interés al que se financia actualmente la empresa:\footnote{
Para incorporar cualquier otro método alternativo puede consultarse [\authorfont{Ver Tema 3.B.3}]
}

\eqblock{
\begin{aligned}
    \boxed{r_{WACC}=i(1-\tau)}
\end{aligned}
}{}

De todas formas, también podemos establecer un umbral mínimo de rentabilidad y fijar un coste de capital endógeno a través de la TIR.\footnote{%
    Se trata de una alternativa habitual en las operaciones. En función de su mandato de inversión, determinan una rentabilidad mínima a obtener en la operación y esta se refleja como el coste del capital. Los fondos de capital de riesgo pueden llegar a demandar tasas de retomo (TIR) de hasta el 25-30\% en las operaciones.
} Por último, para obtener el valor de la empresa debemos realizar la operación en dos pasos, como indicamos antes:

\eqblock{
\begin{aligned}
    &\boxed{\text{Enterprise Value}=\underset{\text{\scriptsize VAN: valor presente del horizontal estimado}}{\sum_{t=0}^T\frac{FCL_t}{(1+r_{WACC})^t}}+\underset{\text{\scriptsize valor terminal descontado al presente}}{\frac{\text{ValorTerminal}_T}{(1+r_{WACC})^T}}}\\[1em]
    &\text{donde},\quad \text{ValorTerminal}_T:\quad \frac{FCL_T(1+g)}{r_{WACC}-g}\qquad \text{o},\qquad \underset{\text{\scriptsize de EBITDA por ejemplo}}{\text{Múltiplos}}
\end{aligned}
}{}

Por tanto, el resultado de la suma de la descomposición del valor nos dará el enterprise value. Como vemos existen diversos métodos para estimar el valor terminal: (i) utilizar la fórmula de una perpetuidad con crecimiento constante;\footnote{%
    Es análogo al modelo de Gordon-Shapiro: un modelo de dividendos crecientes a tasas constantes.

Donde $g$ es la tasa de crecimiento de equilibrio a largo plazo que la empresa puede esperar alcanzar a perpetuidad, teniendo en cuenta tanto la inflación, el crecimiento real como la tasa de reinversión del beneficio operativo. La tasa de crecimiento terminal suele ser inferior a la tasa de crecimiento aplicada en la primera fase, ya que las ventajas derivadas de sinergias, nuevas oportunidades de negocio reducciones de costes son transitorias, pues el entorno competitivo cambia y el sector evoluciona con el tiempo.
} o, (ii) aplicar un múltiplo al que espera que se venda la empresa al final de la primera etapa (más habitual en operaciones de capital riesgo), EBITDA, etc.\footnote{%
    En general, los múltiplos utilizados para estimar el valor residual son convenciones de mercado aplicadas por analistas, banqueros de inversión e inversores de capital riego en función de los distintos sectores
} 

\paragraph{Valoración: Sensibilidad y ajustes}
La ventaja fundamental de este método es que los cambios previstos en los flujos de caja de la empresa objetivo (por ejemplo, las sinergias y cambios en la estructura de costes) pueden modelizarse fácilmente. 

La construcción del modelo proporciona una estimación del valor intrínseco de acuerdo con las hipótesis realizadas, permitiendo analizar la sensibilidad a las mismas (lo cual puede ser un elemento clave en la negociación). También es ventajosa la posibilidad de obtener diversos valores más allá del Valor de Empresa. Por ejemplo, a partir del flujo de caja libre puede obtener el flujo de caja para los accionistas (Free Cash Flow to Equity) y aplicando la tasa de descuento correspondiente obtener el valor del capital propio (Equity Value). 

Es decir, supone un modelo de referencia que se puede adaptar a las necesidades de empresas, analistas e inversores en función de sus objetivos.

Sin embargo, este análisis también tiene una serie de puntos débiles: 
\begin{lnum}
    \item Es dificil de aplicar cuando los flujos de caja libres son negativos en el período de proyección. Por ejemplo, una empresa en rápida expansión puede ser rentable pero tener flujos de caja libres negativos debido a las inversiones necesarias. El valor de empresa procederá principalmente del período posterior ( valor terminal) y será por lo tanto más difícil de estimar;
    \item Estimar los flujos de caja y los beneficios a largo plazo está lejos de ser una ciencia exacta. Existe una gran incertidumbre en las estimaciones para 4-5 años, y una incertidumbre aún mayor a perpetuidad;
    \item Las estimaciones de las tasas de descuento pueden cambiar con el tiempo debido a la evolución de los mercados de capital y/o a cam?ios que afecten a las empresas y al sector en cuestión;
    \item Las previsiones del valor terminal suelen estar sometidas a un gran error de estimación, al verse afectada drásticamente por pequeflos cambios en el crecimiento supuesto y el W ACC.
\end{lnum}

\subsubsection{Valoración por múltiplos}
Un segundo enfoque utilizado para estimar los valores de adquisición es la valoración por múltiplos. Este enfoque a su vez tiene dos métodos: el análisis de compañías comparables y el análisis de transacciones comparables.

\paragraph{Empresas comparables} 
En este enfoque, se debe definir primero un conjunto de otras empresas que son similares a la empresa objetivo a estudiar (frecuentemente denominado Peer Group).  Este grupo puede incluir empresas del sector principal de la empresa objetivo, así como empresas de sectores similares. La muestra debe incluir, en la medida de lo posible, el mayor número de empresas con un tamaño y una estructura de capital similares a los de la empresa objetivo. 

\eqblock{
\begin{aligned}
    \left.\begin{aligned}
        \substack{\text{\scriptsize No intereses}\\\text{\scriptsize No escudo fiscal}}&\underset{\text{óptimo para diferntes apalancamientos}}{\left\{\begin{aligned}
            &\text{EV/EBITDA}\\[0.5em]
            &\text{EV/EBITDA}\\[0.5em]
            &\text{EV/Ventas}\\[0.5em]
            &\text{EV/FCF}
        \end{aligned}\right.}\\[0.5em]
        &\left\{\begin{aligned}
            &\text{PER}\\[0.5em]
            &\text{EquityValue/FCLE}\\[0.5em]
            &\text{Precio acción/Ventas}\\[0.5em]
            &\text{EquityValue/ValorContable}\\[0.5em]    
        \end{aligned}\right.
    \end{aligned}\right\}\underset{\text{\scriptsize ajustes}}{\Longrightarrow}
    \underset{\Longrightarrow \text{DEA}}{\left\{\begin{aligned}
        &\text{Tamaño}\\[0.5em]
        &\text{Mercados geográficos}\\[0.5em]
        &\text{Empleados}\\[0.5em]
        &\text{Ventas}\\[0.5em]
    \end{aligned}\right.}
\end{aligned}
}{}

Como vemos, la elección de los comparables es un momento crucial. Si no disponemos de una muestra suficiente, es posible llevar a cabo ajustes en base a criterios objetivos que determinan la distancia relativa de la empresa objetivo con aquellas de nuestro grupo control. Por ejemplo, mediante técnicas estocásticas de análisis de fronteras estocásticas (DEA, envolvente de datos).\footnote{%
    Las condiciones de los mercados de capitales y la disponibilidad de empresas comparables dificultan este tipo de análisis por lo que es habitual recurrir a rangos de valoración. Los analistas de una operación corporativa suelen examinar la media, la mediana y el rango de cualquiera de los múltiplos elegidos y, a continuación, aplican esos valores a las estimaciones correspondientes de la empresa objetivo para elaborar una estimación del valor de la empresa. 

Es también frecuente que cada múltiplo (EV/EBITDA, P/E, etc.) produzca una estimación diferente del valor. Lo ideal es que estos valores tiendan a converger, aumentando la confianza en el conjunto de la estimación. Sin embargo, en la medida en que esto no sucede y los valores tienden a ser divergentes, se deben aplicar juicio y experiencia para decidir qué múltiplos o medidas de valor relativo producen los valores más precisos. En ocasiones se recurre a la construcción de múltiplos sintéticos, realizando una media ponderada o aritmética de varios múltiplos.
} Sea como fuere, una vez definido el universo de empresas comparables, el siguiente paso consiste en calcular los múltiplos (medidas de valor relativo).

\paragraph{Transacciones comparables} 
Otro enfoque estrechamente relacionado con el de empresas comparables es utilizar el detalle de transacciones recientes de adquisición para hacer estimaciones directas del valor de adquisición de la empresa objetivo.

El primer paso, análogo al proceso anterior, es la recopilación de una muestra relevante de operaciones recientes. Para ello es habitual recurrir a bases de datos especificas en las que se busca un equilibrio entre la confidencialidad de las operaciones y el valor como bien público que esta información aporta al conjunto del sector empresarial. La muestra debe ser lo más amplia posible, pero limitada a empresas del mismo sector, o al menos estrechamente relacionadas. Una vez identificadas las transacciones, se examinan los mismos múltiplos que en el método anterior. 

En este caso comparamos los múltiplos realmente pagados por empresas similares en otras operaciones de fusión y adquisición, en lugar de los múltiplos \textit{cotizados}. Al igual que antes, los analistas suelen utilizar la media, la mediana y el rango de los múltiplos, y aplican criterio y experiencia al aplicar esa información para estimar el valor. 

\paragraph{Ventajas}
En general, estos métodos proporcionan una aproximación razonable del valor de una empresa objetivo en relación con empresas similares del mercado. 

Se parte de la base de que activos similares deben valorarse de forma parecida en el mercado. La principal ventaja es que la mayoría de los datos necesarios están fácilmente disponibles y las estimaciones de valor se derivan directamente del mercado, evitando que las hipótesis y supuestos estén lejos de lo que el mercado está valorando. 

No obstante, en la práctica puede ser un reto importante encontrar empresas o transacciones similares y comparables. En estos casos se debe aceptar que las empresas son diferentes y tratar de controlar explícita o implícitamente esas diferencias en cuanto a crecimiento, riesgo y medición de los flujos de caja. Para ello se utilizan desde variables estándar como medias o medianas del sector hasta metodologías más sofisticadas con modelos de regresión multivariante en los que se identifican y controlan las variables relevantes. La sensibilidad al mercado es un punto fuerte de este método pero también puede ser considerado una desventaja, al estar sujeto a los ciclos de mercado y las posibles valoraciones erróneas del mismo.

Además, otro punto débil muy relevante es la dificultad de adaptar la valoración a la transacción en cuestión. La valoración a través de múltiplos produce una estimación aproximada sobre dónde debería cotizar la empresa objetivo como acción en el mercado, en relación con las empresas de la muestra. No obstante, este valor no tiene por qué coincidir con el valor de adquisición a la hora de valorar una operación de fusión o adquisición. El valor de adquisición debe tener en cuenta las sinergias, reducciones de coste y otros elementos vinculados al valor que se espera crear con la operación. Debe sumarse también una prima de adquisición. Es decir, la cantidad en la que el precio de adquisición de cada acción debe superar el precio actual de las acciones para que los accionistas actuales estén dispuestos a ceder el control de la empresa a un adquirente. Para incorporar los aspectos específicos de la operación y calcular la prima de adquisición pertinente para una operación, los analistas suelen recopilar una lista de las primas pagadas por empresas similares. Preferiblemente, los cálculos serán del pasado reciente porque los valores de adquisición y las primas tienden a variar con el tiempo y los ciclos económicos. Esto nos lleva al segundo método de valoración por múltiplos: el análisis de transacciones comparables.

Frente a la valoración de empresas comparables, el punto fuerte de las transacciones comparables es que no es necesario estimar por separado una prima de adquisición, al derivarse directamente de las transacciones comparables. El uso de precios establecidos a través de otras transacciones recientes reduce el riesgo de litigios para los consejos de administración y los directivos de ambas empresas en relación con el precio de la transacción. El supuesto implícito de este método es que el mercado de fusiones y adquisiciones ha determinado correctamente el valor intrínseco de las empresas objetivo. No se debe obviar el impacto que las condiciones macroeconómicas y financieras han podido tener en estas transacciones, o el impacto de los cambios en la estructura del sector.

Asimismo, es posible que no haya transacciones comparables, o que no haya un número suficiente de ellas. En estos casos, los analistas pueden intentar utilizar como second best los datos de sectores similares.

\subsubsection{Opciones Reales}
[...]

Así, el valor de una opción viene determinado esencialmente por seis variables\footnote{Ampliamos la formula, siguiendo a Damodaran, para incluir D.
}

\begin{center}
\begin{tabular}{|p{0.33\linewidth}|p{0.33\linewidth}|p{0.33\linewidth}|}
\hline
\textbf{OPCIÓN REAL} & \textbf{VARIABLE} & \textbf{OPCIÓN FINANCIERA} \\
\hline
Valor activos a adquirir & $S$ (spot) & Precio subyacente \\
\hline
Inversión requerida & $X$ (strike) & Precio ejercicio \\
\hline
Aplazamiento & $t$ & Precio ejercicio \\
\hline
Riesgo de la inversión & $\sigma^2$ & Varianza rendimiento de activos financieros \\
\hline
Tipo de interés (libre de riesgo) & $\gamma$ & Tipo de interés (libre de riesgo) \\
\hline
Flujos de caja perdidos en caso de no ejercer la opción & $D$ & Dividendos del activo subyacente \\
\hline
\end{tabular}
\end{center}

En el caso de opciones reales: • S: Valor actual esperado de los flujos de caja del activo o inversión • X: Precio a pagar por el activo real o inversión a realizar • t: Tiempo durante el cual se puede demorar la decisión de inversión • o-2 : Volatilidad de los flujos de caja del activo real/inversión • r: Tipo de interés libre de riesgo (valor temporal del dinero) • D: flujos de caja a los que se renuncia mientras no se ejerza la opción Como se observa en la propia fórmula, el modelo Black-Scholes asume normalidad de las distribuciones de flujos de caja, un supuesto que no se cumple en la práctica. Además, se trata de un modelo rígido en lo que refiere al ejercicio temprano de la opción. Así, dadas las características de estas opciones reales en el contexto corporativo, el método binomial es en la práctica el más utilizado.A pesar de las analogías entre la valoración de opciones financieras y la de opciones reales, estaspresentan una serie de limitaciones: • Mientras que en el caso de una opción financiera, el strike o precio de ejercicio (X) es fijo, en una opción real, la inversión o desembolso necesario es una variable que puede ser volátil en función de la evolución de los precios de los proveedores, la disponibilidad de mano de obra, etc. • Del mismo modo, S, el valor actual solo se conoce de manera aproximada, al no existir necesariamente un valor cotizado observable en todo momento. • La incertidumbre o volatilidad de los flujos de caja es externa en el caso de una opción financiera, pero en una opción real existe la posibilidad de influir sobre la misma (mitigando o no riesgos como el de tipo de cambio, influyendo sobre el comportamiento de la competencia, etc.)En definitiva, la valoración de opciones reales es más compleja al tratar de modelizar un mundo real en que las contingencias y eventos posibles son más amplios que en el caso de los activos financieros. Al igual que en otras metodologías, el juicio a aplicar por parte de los analistas tiene un papel clave.

\subsubsection{Suma de partes: Conglomerados}
El método de suma de las partes (SOTP, por sus siglas en inglés) consiste en valorar y sumar los diferentes activos, divisiones o filiales de la empresa. 

Es un método muy adecuado para grupos o conglomerados diversificados, para los que las proyecciones a nivel consolidado dan una visión demasiado global. No se trata como tal de una metodología diferente a las que se han visto con anterioridad. Consiste en analizar de manera independiente el valor de los activos y pasivos de cada línea de negocio, descomponiendo la empresa en diferentes partes o componentes y estimar el valor de cada una de ellas por separado. Posteriormente, se suman estos valores individuales para obtener el valor total de la empresa. 

El proceso general implica los siguientes pasos:
\begin{lnum}
    \item Identificación de los componentes: Se analiza la empresa para identificar los diferentes componentes que la confonnan. Estos componentes pueden incluir divisiones de negocios, unidades operativas o activos específicos que tienen un impacto significativo en el valor de la empresa;
    \item Análisis individual de cada componente: Se realiza un análisis detallado de cada componente identificado. Esto implica evaluar sus flujos de efectivo futuros, riesgos asociados, crecimiento esperado y otros factores relevantes para detenninar su valor individual;
    \item Valoración de cada componente: bien a través de un modelo de descuento de flujos de caj a o a través de métodos de valor relativo (múltiplos). Como en cualquier otra valoración, se deben tener en cuenta las tendencias del mercado, las condiciones económicas y otros factores relevantes y específicos (por ejemplo, tasas de descuento diferentes);
    \item Suma de los valores individuales: Una vez que se han estimado los valores individuales de cada componente mediante el descuento de los flujos de efectivo futuros, se suman para obtener el valor total de la empresa;
    \item 
\end{lnum}

Una de las principales ventajas de este método es comprobar que en muchas ocasiones la valoración de los conglomerados está por debajo de la suma de sus partes. Esto puede ser fruto de un problema de información incompleta y/o asimétrica (el mercado no percibe correctamente el valor total de la cómpañía); o una señal de que la dirección de la empresa no está maximizando su valor. Estas divergencias en la valoración pueden llevar a procesos de desinversión con el fin de maximizar el valor para los accionistas.


\subsection{Empresas cotizadas: Capitalización bursátil}
Hasta aquí hemos visto las diferentes opciones metodológicas que tenemos para evaluar el enterprise value de una empresa no cotizadas. Esto nos permite, como ya hemos visto, alcanzar el equity value restando al enterprise value el valor correspondiente a los acreedores financieros.

En el caso de empresas cotizadas, el criterio de valoración es radicalmente diferente. ¿Por qué? Porque el valor de la empresa que le corresponde a los accionistas viene determinado por los mercados financieros. De esta forma, para observar este valor, analizar la rentabilidad de la operación y, en caso de que esta sea rentable, realizar la Oferta Pública de Adquisición, deberemos fijarnos en la \textbf{capitalización bursátil} de la empresa.

\subsubsection{Ecuación de canje: Valor relativo}
La magnitud de las empresas cotizadas conlleva que la mayoría de estas operaciones se instrumentalicen a través de una fusión o canje de acciones. 

En este contexto, es necesario determinar la \textbf{ecuación de canje}, es decir, la relación entre el número de acciones a emitir por cada acción de la otra empresa que se va a recibir. Esta ecuación de canje determinará a su vez el valor relativo de una compañía frente a otra, fruto del ejercicio de valoración realizado y del proceso de negociación de la operación. La ecuación de canje y el valor relativo están relacionados por la siguiente fónnula: 

\eqblock{
\begin{aligned}
    \underset{\text{\scriptsize(Nº de accioens $A$ x 1 acción $B$)}}{\text{Ecuación de canje}}=\underset{\text{\scriptsize ($B$ vale $x$ veces $A$: bursátil)}}{\text{Ratio valor relativo}}\cdot\frac{\text{Nº acciones $A$ (pre-fusión)}}{\text{Nº acciones $B$ (pre-fusión)}}
\end{aligned}
}{}

Como se puede observar claramente, la ecuación de canje nos dará directamente el valor de mercado de las acciones de $A$ por cada acción de $B$. No obstante, cabe suponer que la ecuación de canje finalmente pactada (u ofertada) será mayor por la prima de control exigible por la empresa objetivo. De esta forma, podremos hallar la valoración relativa finalmente acordada a raíz de la ecuación de canje pactada.

\clearpage
\section{Evidencia Empírica: Implicaciones de Política Económica}

Hasta ahora hemos visto las operaciones de M\&A desde una óptica empresarial. Hemos visto como se clasifican, tanto jurídicamente como económicamente, y hemos visto diferentes criterios de valoración de los que se pueden servir los generentes y analistas en los procesos. Ahora bien, no hemos abordado un campo crucial: ¿qué determina el crecimiento? ¿qué implicaciones en términos de Política Económica?

Estas dos cuestion son profundamente relevantes cuando tratamos los procesos de fusión y adquisición desde una perspectiva de supervisión y control de la competencia y, como tal, son objeto de otras partes del programa. No obstante, sí cabe analizar de forma indirecta estas implicaciones a través del análisis empírico de estas operaciones y, en consecuencia, la búsqueda de los determinantes del crecimiento inorgánico. Pasamos a analizar con más detalles los principales motivos que guían las operaciones de fusiones y adquisiciones desde un punto de vista práctico y con un enfoque principalmente microeconómico.


Cualquier proceso de fusión acarrea reducciones de costes que, comportan a menudo la eliminación de puestos de trabajo y el cierre de instalaciones productivas y oficinas comerciales, representan una piedra de toque neurálgica en pos del aumento de la rentabilidad de los accionistas. Tras esos empleos que desaparecen, detrás de esos cierres de plantas y factorías, hay todo un mundo, más grande o más pequeño, de industria auxiliar, de proveedores de bienes y servicios, de suministradores, que desaparece. Los centros de decisión empresarial se trasladan, se mueven.


\subsection{Evolución de las actividad de M\&A: Ciclos}
El análisis empírico de la actividad de fusiones y adquisiciones lleva a observar que esta se produce a través de oleadas, o ciclos, en el número de operaciones y su volumen, que sucede a todos los niveles (local, sectorial y global). 


\imagenfit{Historico_MA.png}{Number of M\&A Operations in the U.S. from 1897 to the present}{Gaughan(1999), Nelson (1959), Historical Statistics of the U.S. – Colonial Times to 1970, Mergerstat Review.}


Las fusiones no son una moda pasajera que haya irrumpido con fuerza en la encrucijada entre siglos que vivimos, ni una ola que viene y va. Han calado hondo, alterando sustancialmente las coordenadas del panorama empresarial a lo largo y ancho del vasto mapamundi. Es cierto que algunas de esas fusiones responden a situaciones puramente coyunturales pero muchas otras han representado la vía para prosperar manteniendo unas determinadas posiciones hegemónicas. Posiblemente, haya algunas operaciones de concentración empresarial que sorprenden a propios y extraños como, en su momento, dentro del año 2002, fue la consolidación del proceso de integración entre Hewlett- Packard y Compaq que, aún hoy, transcurrido un lapso de tiempo prudencial sigue generando sus incógnitas y determinados escepticismos.



Lo que enfatiza que los distintos determinantes del crecimiento inorgánico se interrelacionan entre sí. 


Los factores macroeconómicos y macrofinancieros (por ejemplo, los tipos de interés) pueden tener un impacto tan relevante como los microeconómicos. En los ciclos de fusiones y adquisiciones, el número de operaciones en su pico puede ser varias veces superior al que se produce en su punto más bajo. 

Es habitual observar como en la fase alcista del ciclo se producen cada vez más operaciones y a precios/múltiplos cada vez más elevados, hasta que algún evento rompe la tendencia y paraliza la actividad por completo. 


Existen dos líneas teóricas principales que tratan de explicar por qué se producen estas oleadas: 
En primer lugar, la teoría que podemos considerar tiene un corte más neoclásico, se centra en aspectos microeconómicos y postula que las oleadas de fusiones ocurren cuando las empresas en industrias específicas reaccionan a shocks económico (regulación/desregulación, aparición de nuevas tecnologías o productos y/o servicios sustitutivos), lo cual explica por qué la actividad de fusiones se agrupa por industrias. Eltamaño y la duración de cada oleada dependen en gran medida del número de industriasinfluenciadas durante un mismo espacio temporal. 

La segunda línea de desarrollos teóricos sugiere que las valoraciones del mercado provocan las oleadas de fusiones. Cuando las valoraciones de las empresas se desvían de los fundamentales, los directivos tienden a utilizar las acciones sobrevaloradas de sus empresas como instrumento para comprar activos de empresas infravaloradas (o menos sobrevaloradas), lo cual explica la correlación entre la actividad de fusiones y el rendimiento de los mercados de acciones. En consecuencia, la teoría de la sobrevaloración postula que habrá más adquisiciones en períodos de burbujas o de altas valoraciones bursátiles, lo cual puede observarse empíricamente. 





\subsection{Determinantes del crecimiento inórganico: Microeconómicos (TOI)}
La Teoría de la Organización Industrial ha estudiado el crecimiento de las empresas principalmente desde el punto de vista de su tamaño eficiente

Los factores microeconomicos son un motivo básico, condicionados por el entorno macroeconómico e institucional. A todo ello hay que sumar las razones vinculadas a los incentivos y el comportamiento de los agentes involucrados (directivos, empleados, intermediarios, gobiernos ... ). Las desarrollamos a continuación: 

\subsubsection{Sinergias: Impacto Microeconómico}
Un determinante fundamental desde el punto de vista microeconómico son las sinergias que pueda generar una operación. En resumidas cuentas: \textbf{que la suma del conjunt o sea superior al valor de las partes por separado}.

Estas sinergias se producir vía múltiples canales, mejorasen ingresos, mejoras en costes, mejoras de gestión, financieras y fiscales o incluso vía diversificación de riesgo.\footnote{%
    Por un lado, las economías de escala y alcance (integración horizon tal) permiten una disminución de los costes medios al aumentar el volumen de producción y/o la gama de productos. En general se dan en operaciones de tipo horizontal, es decir, en empresas con parecido ámbito de actividad, aunque también en conglomerados. En los casos de economías de integración vertical, se justifican las operaciones de crecimiento externo vertical para aproximarse más con sus productos al consumidor final (integración hacia adelante) o a la fuente de materias primas (integración hacia atrás). 

Por el lado de los ingresos, la fuente de sinergia se produce cuando la entidad fusionada alcanza mayores ventas y crecimiento que el que las dos empresas independientes podrían alcanzar. Esta mejora puede proceder de una ofe rta de productos racionalizada (productos complementarios) o del mejor aprovechamiento de las redes de distribución comercial. 

En ocasiones se producen sustanciales mejoras de gestión o gobierno corporativo. En cuyo caso la operación es un mecanismo para remplazar a un equipo directivo por otro. Aunque no es el único método para mejorar la dirección de una compañía, es práctico y simple, superando resistencias al cambio orgánico. Los directivos nunca tienen un incentivo a prescindir de ellos mismos y en empresas cotizadas muchos accionistas ejercen una influencia poco directa sobre la dirección. En definitiva, estas operaciones pueden servir de mecanismo de disciplina de mercado.
} Por ello, la medición de impacto es tremendamente arbitraria pero, en esencia consiste en la diferencia entre el valor de la empresa combinada y los valores individuales de las empresas:\footnote{%
    Dado que las empresas adquirentes deben pagar generalmente una prima por hacerse con el control de la firma adquirida, es importante calibrar que esta prima está suficientemente justificada por el potencial valor de la sinergia. Un método habitual es utilizar un modelo de descuento de flujos de caja adaptado: por un lado, se debe tener en cuenta los flujos incrementales y por otro lado, utilizar como tasa de descuento el coste de los recursos propios de la empresa adquirida (en lugar del coste de los recursos de la adquirente). Esta tasa de descuento estaría reflejando el riesgo de los flujos de caja incrementales.
}

\eqblock{
\begin{aligned}
    \underset{\text{\scriptsize VA de la operación}}{\text{Sinergías}}=V(\text{adquirente})-V(\text{objetivo})
\end{aligned}
}{}

Empíricamente, los investigadores se basan en datos del mercado de valores o en datos contables para inferir las ganancias de sinergia reales que se acaban materializando. En estudios más recientes se ha priorizado la medida contable de mejoras de eficiencia: el cambio en elrendimiento operativo anormal de la empresa combinada en comparación con el rendimiento operativo anormal previo a la fusión de la empresa compradora y la empresa objetivo.\footnote{%
    WANG. C. y F. XIE (2009): \textit{Corporate governance transfer and synergistic gains from mergers and acquisitions}. Review of Financia! Studies, 22, 829-858.

Por lo general, el rendimiento operativo se define como el retomo sobre los activos totales (ROA) y el rendimiento anormal se mide en comparación con la mediana de la industria o una empresa de control elegida en función de la industria, el tamaño y el nivel de rendimiento operativo previo al evento corporativo. Por lo general, se utiliza el promedio del rendimiento operativo anormal de tres años después de la adquisición y se compara con el rendimiento operativo anormal en el período previo a la adquisición.
}

De forma no exhaustiva, algunas fuentes de ganancias por sinergias son las siguientes.

\paragraph{Ventajas fiscales}
Por otro lado, se añade un componente fiscal cuando una empresa rentable puede adquirir a una que no lo es para, entre otros aspectos, disminuir su carga fiscal. 

En estos casos la sinergia procede de la plena utilización de escudos fiscales ya que es habitual que haya escudos fiscales infrautilizados si la compañía no explota su capacidad de endeudamiento o si la empresa adquirida tiene bases imponibles negativas pendientes de compensar. Una adquisición también permite aflorar un valor intangible que no estaba reflejado en los estados financieros, dando lugar a un fondo de comercio (que aumentará los gastos por depreciación). 

No obstante, la compensación de bases imponibles negativas se ha limitado en los últimos años en algunas jurisdicciones, disminuyendo el incentivo a este tipo de operaciones.

\paragraph{Diversificación del riesgo}
Otro motivo menos visible es la diversificación del riesgo, aunque es un motivo habitual en las fusiones y adquisiciones. No obstante, se trata de uno de los motivos más discutibles, en la medida que la literatura económica apunta a que la diversificación en sí misma dentro de la empresa (a nivel geográfico e industrial) puede no aportar un valor adicional frente al que el inversor individual puede conseguir en su cartera de acciones.\footnote{%
    S.B. MOELLER and F. P. SCHLINGEMANN. \textit{Global Diversification and Bidder Gams: A Companson between Cross-Border and Domestic Acquisillons.} Journal of Banking and Finance 29. no. 3 (March 2005): 533-564.

Esta evidencia habría que contraponerla con los supuestos detrás del valor que aporta la diversificación individual para el inversor. En EE.UU los mercados de capitales (y en particular, el de renta variable) son profundos, pero puede no ser el caso en otros mercados (como en España, donde el inversor tiene limitadas sus opciones para diversificar entre industrias).
} 

\eqblock{
\begin{aligned}
    &\downarrow \beta_A\,\,\Longrightarrow\,\, \downarrow E_t[r_{WACC}]=\gamma+\underbrace{\beta_A(\underbrace{\mathbb E_t[r_M]-\gamma}_{=\lambda_M})}_{=\lambda}\\[1em]
    &
\end{aligned}
}{Sinergias: Diversificación del riesgo y coste de capital (CAPM). }

A pesar de esto, algunos autores consideran que la diversificación puede ser efectivamente una fuente de sinergias: pueden crear mercados de capitales internos, que permiten asignar recursos entre divisiones sin las fricciones e ineficiencias de los externos.\footnote{%
    Las operaciones transfronterizas también pueden servir como un medio para ampliar el perímetro de la empresa adquirente de forma que esta pueda internalizar los beneficios que se perderían debido a diversas fricciones del mercado, siendo esto más relevantes en el ámbito internacional.

J. DOUKAS y N.G. TRAVLOS (1988), \textit{The effect of corporate multinationalism on on shareholders' wealth: Evidence from international acquisitions}, Journal of Finance, 43(5), 1161-1175.
} 

\subsubsection{Incentivos: Motivos de agencia}
Por otro lado, algunas tesis apuntan a que los directivos pueden llevar a cabo adquisiciones en contra de los intereses de los accionistas en persecución de sus propios objetivos. Las causas de estos problemas de agencia son variadas, entre las que encontramos:
\begin{lnum}
    \item \textbf{Suavización de beneficios}. Fusiones para diversificar las actividades de la empresa y suavizar la volatilidad de los beneficios, maximizando la probabilidad de mantener sus empleos en el tiempo. Esto iría en contra de los intereses de los accionistas, que serían capaces de diversificar sus inversiones a un menor coste;
    \item \textbf{Incentivos: Valor de la empresa}. Dado que la remuneración y prestigio de los equipos directivos está muy relacionada con el tamaño de la empresa, estos pueden tener un incentivo a maximizar el tamaño de la empresa (no necesariamente el valor de los accionistas). Una forma de maximizar este tamaño es participando en fusiones y adquisiciones;
    \item \textbf{Excedentes de tesorería}. En lugar de redistribuir el efectivo acumulado y reducir el tamaño de la empresa, los directivos pueden tender a involucrarse proyectos personales y adquisiciones no rentables o con una creación de valor neta negativa. Los costos de agencia asociados a estos comportamientos podrían ser una fuente de destrucción de valor (frente a las sinergias).
\end{lnum}

En esta línea y en cuanto a la alineación de incentivos, varios estudios sugieren que los directivos de la empresa adquirente cuya riqueza personal está más estrechamente vinculada al valor de la empresa (mayor participación en el accionariado), tienden a tomar mejores decisiones de adquisición. Otros análisis apuntan a la importancia de los excedentes de tesorería (flujos de caja por encima de lo que la empresa puede invertir de manera rentable): las empresas adquirentes que disponen de mucho efectivo (por encima de lo normal en una industria) serían más propensas a intentar adquisiciones, en particular operaciones de diversificación, y en términos agregados tienden a destruir valor para el accionista y a sufrir descensos anormales en el rendimiento operativo.


\subsubsection{Crecimiento: Poder de mercado}
No obstante, uno de los determinantes que no cabe perder nunca de vista y ha venido siendo contrastado empíricamente los últimos años es el de impulsar el crecimiento de la empresa y conseguir mayor poder de mercado. 

\textbf{WARREN BUFFET}

El crecimiento interno puede ser un proceso lento e incierto y el crecimiento externo se caracteriza principalmente por ser mucho más rápido. Esto puede llevar a priorizar el crecimiento a través de las operaciones corporativas para ganar tamaño más rápido y/o modificar el perímetro de la compañía. En muchos sectores se observa que la distribución de diversas ratios indicativas del poder de mercado están adoptando distribuciones cada vez más asimétricas desde los años 80's (en la ratio beneficios/ventas, la ratio salario/gastos, etc.). Esto es indicativo de un fenómeno especialmente importante: cada vez hay más desigualdad intrasectorial e intersectorial. Son muchas las causas que han motivado estos cambios, pero lo que realmente preocupan son las consecuencias que puedan acerrear (se puede seguir aquí).

\paragraph{Capacidades y recursos únicos - Activos específicos} 
Muchas empresas emprenden una fusión o una adquisición para obtener ventajas competitivas o para reforzar los recursos que les faltan. Cuando una empresa no puede crear internamente de forma rentable las capacidades necesarias para mantener su éxito futuro, puede tratar de obtenerlas a través de una operación de crecimiento externo. El objetivo pueden ser competencias o recursos específicos de los que carece, como un departamento de investigación con éxito, una red de distribución comercial, o capital humano. De la misma fonna, puede haber una complementariedad entre empresas pequefias y grandes. Muchas pequeñas empresas son adquiridas por grandes empresas porque éstas les pueden aportar elementos que son necesarios para el éxito de un producto. A la empresa compradora le puede permitir entrar en nuevo nicho rentable y a la adquirida obtener el capital necesario para el desarrollo del producto. Esto algo común en sectores de alto contenido tecnológico corno el farmacéutico.


\subsection{Determinantes del crecimiento: Macroeconómicos}
Al estudiar los ciclos de M\&A, se ha obersado que suelen coincidir con mercados de capitales en alza y una recuperación de la actividad económica. La siguiente correlación de la actividad de M\&A con una serie de factores relacionados con los mercados de renta variable y de bonos muestra esta tendencia:

\imagenfit{OFI_REPORT.png}{(Izq.) Actividad medida como \% de la capitalización bursátil total del mercado EE.UU | (Dcha.) Correlación entre facilidad de crédito y la variación anual entre actividades de M\&A}{\href{https://www.caia.org/sites/default/files/AIAR_Q3_2015-07_M\%26A_CretinDieudonneBoucha.pdf}{\textit{M\&A Activity: Where Are We In the Cycle?}. CRETIN, DIEUDONNÉ y BOUACHA (2015). OFI Asset Management}}

El primer factor es la capitalización total del mercado estadounidense. Para cada trimestre del período analizado, se compara el valor total de las operaciones anunciadas en los últimos 12 meses móviles con la capitalización bursátil total del mercado de EE. UU. y se muestra que este ratio se sitúa en un rango del 2\% al 4\% en los mínimos del ciclo y entre el 9\% y el 12\% en los picos del ciclo. El segundo factor es la encuesta Senior Loan Officer de la Reserva Federal, que mide las condiciones de crédito para empresas y consumidores estadounidenses. Un nivel alto refleja condiciones de financiación difíciles y un nivel bajo refleja condiciones favorables. 

\imagenfit{OFI_REPORT_2.png}{(Izq.) Actividad medida como \% de la capitalización bursátil total del mercado EE.UU | (Dcha.) Número de operaciones canceladas en relación con las operaciones completadas}{\href{https://www.caia.org/sites/default/files/AIAR_Q3_2015-07_M\%26A_CretinDieudonneBoucha.pdf}{\textit{M\&A Activity: Where Are We In the Cycle?}. CRETIN, DIEUDONNÉ y BOUACHA (2015). OFI Asset Management}}

El tercer factor es un indicador que mide el número de mejoras de oferta esperadas por el mercado en operaciones de M\&A sobre compañías cotizadas. En ciertos casos, cuando se anuncia una operación, el precio de la acción de la empresa objetivo puede situarse por encima del precio ofrecido por el comprador potencial. Esto indica que los inversores esperan una mejora de la oferta por parte del comprador o una contraoferta de un tercero. Se observa claramente que los máximos y mínimos de este indicador coinciden con las variaciones del ciclo de actividad de M\&A. El cuarto factor se refiere a operaciones de M\&A en empresas cotizadas. Se trata de la tasa de fracaso de operaciones calculada sobre 12 meses móviles, es decir, el número de operaciones canceladas en relación con las operaciones completadas. Diversos factores pueden provocar el fracaso de una operación: oposición de los accionistas de la empresa objetivo o de las autoridades de competencia, problemas de financiación, deterioro de las condiciones de mercado, oposición de los accionistas del comprador, etc. La tasa media de fracaso es aproximadamente del 7\%. Un aumento significativo por encima de este nivel refleja un incremento general de la aversión al riesgo en el entorno de M\&A. Una vez más, los periodos de descenso en la actividad de M\&A coinciden con momentos en los que la tasa de fracaso es elevada.



\clearpage
\section{Conclusión} 


\subsection{Recapitulación}


\subsection{Extensiones y relación con otras partes del Temario}



\subsection{Opinión}

\subsubsection{Idea final (Salida o cierra)}

\clearpage
\pagenumbering{gobble}
\MyBibliography



\MyBibItem{DAWID y GATTI}{Chapter 2 - Agent-Based Macroeconomics}{2018}{https://doi.org/10.1016/bs.hescom.2018.02.006}


% ANEXOS
\clearpage
\anexo{\textbf{Preguntas Test}}
%\input{\TemaCode/questions.tex} 



\end{document}